\section{Claudeのテキスト・ウォーターマーク議論が継続}
AnthropicがClaudeの出力に適用する不可視のテキスト・ウォーターマークをめぐる議論が引き続き活発だった。投稿では、語の選択に統計的なパターンを持たせて検出可能性を確保する仕組みや、EU AI Act対応との位置づけが紹介された。一方、コードやSEO、業務文書でAI支援を利用する場合の扱い、検出器へのアクセス、誤検出への懸念なども共有されている。機能自体は観測窓より前から存在し、今回はその影響をめぐる反応が中心である。
\dailyxsource{議論投稿例}{https://x.com/newslit/status/2089462783528120515}
\dailyxnote{機能自体の導入は観測窓以前。本文は観測窓内で継続した議論を扱う。}
